We coded the adjudication design of 22 published studies that relabel, audit or clean machine-learning benchmarks (Table~\ref{tab:literature}), each code backed by a verbatim methods quote checked mechanically against the paper's full text; the protocol is in the Online Appendix. The categories are DR (only items flagged by a model or algorithm are adjudicated), DR-multi (the flag is the union over a panel), CRS (flagged items plus a sample of unflagged ones), RANDOM (a sample drawn independently of models), FULL (every item) and OTHER (hybrids).

Eight studies use DR or DR-multi, and five of the six hybrids contain a step in which humans see only items selected by model outputs (one of them, SWE-Bench+, reviews only items the model solved; \citealp{aleithan2024swebench}). In all, 13 of the 22 studies adjudicate only model-selected items for at least part of the work. Of the eight DR studies, seven re-report accuracies or rankings on the corrected labels, seven evaluate a flagging model afterwards, and only four acknowledge that unflagged items can be mislabelled. \citet[p.~4]{beyer2020are} state the opposite assumption outright: ``In the event that all models agree with the original ImageNet label, we can safely retain the original without re-evaluation.''

Several studies contain evidence against that assumption. A random unflagged ImageNet image per class raised the estimated error rate of \citet{northcutt2021pervasive} from 6\% to about 20\%, but the estimate was not used to correct labels or rankings. A full review of SciFact found 3 of its 11 label errors among items the automated screen had not flagged \citep{sylvestre2026gold}. \citet{vendrow2025do} call their error counts lower bounds, since ``our approach could miss any bad question for which all LLM solutions happen to agree with the benchmark solution'', and \citet{chong2022detecting} warn that ``the cleaning process may bias any such benchmarks toward the models most similar to the model used to clean the data.'' The design closest to a composite reference standard \citep{alt2020tacred} finds an 8.9\% revision rate in a control group of items most models got right, but does not propagate it to unsampled items. Two studies re-annotate random samples, including MMLU-Redux \citep{gema2025are}, our ground truth in Section~\ref{sec:empirical}.

Sections~\ref{sec:bias}--\ref{sec:twophase} make these remarks quantitative. For single-flagger designs, Corollary~\ref{cor:advantage} implies that every reported improvement of another model over the flagger is a lower bound and every reported advantage of the flagger is suspect. Other flagging rules, item removal and multi-label scoring fall outside the corollary; for them Proposition~\ref{prop:bias} still applies when relabelled items are scored against the adjudicated label.

\begin{table}[t]
\centering
\scriptsize
\begin{tabular}{llllll}
\toprule
Study & Benchmark & Design & Flagger & Re-reports & Notes unflagged \\
 & & & evaluated & acc./ranking & may be wrong \\
\midrule
\citet{reiss2020identifying} & CoNLL-2003 & DR & yes & yes & no \\
\citet{northcutt2021pervasive} & 10 test sets (ImageNet, CIFAR, \dots) & DR & yes & yes & yes \\
\citet{vasudevan2022when} & ImageNet (multi-label) & DR & yes & yes & yes \\
\citet{chong2022detecting} & IMDB, Amazon, \dots & DR & partial & partial & yes \\
\citet{nahum2025are} & TRUE, SummEval & DR & yes & yes & no \\
\citet{ansari2026how} & 4 physics benchmarks & DR & yes & yes & no \\
\citet{beyer2020are} & ImageNet (ReaL) & DR-multi & yes & yes & no \\
\citet{vendrow2025do} & 15 subsets (GSM8K, SQuAD2.0, \dots) & DR-multi & yes & yes & yes \\
\citet{alt2020tacred} & TACRED & CRS & yes & yes & yes \\
\citet{shankar2020evaluating} & ImageNet & OTHER & yes & yes & yes \\
\citet{rucker2023cleanconll} & CoNLL-2003 & OTHER & yes & yes & yes \\
\citet{wang2024mmlupro} & MMLU-Pro & OTHER & yes & no & no \\
\citet{aleithan2024swebench} & SWE-bench & OTHER & yes & yes & no \\
\citet{ye2026scalable} & MedCalc-Bench & OTHER & no & yes & yes \\
\citet{land2026auditing} & 32 subsets (RewardBench 2, \dots) & OTHER & -- & no & yes \\
\citet{chowdhury2024introducing} & SWE-bench (Verified) & RANDOM & -- & yes & no \\
\citet{gema2025are} & MMLU (Redux) & RANDOM & -- & yes & no \\
\citet{wang2019crossweigh} & CoNLL-2003 & FULL & -- & yes & -- \\
\citet{stoica2021retacred} & TACRED & FULL & -- & yes & -- \\
\citet{zhai2026hleverified} & Humanity's Last Exam & FULL & -- & yes & -- \\
\citet{brunello2026fixing} & FOLIO, MALLS & FULL & partial & yes & -- \\
\citet{sylvestre2026gold} & SciFact & FULL & yes & yes & -- \\
\bottomrule
\end{tabular}

\caption{Adjudication designs of published benchmark relabelling and audit studies, coded from verbatim methods quotes (protocol in the Online Appendix). ``Flagger evaluated'' means that a model used to select items for adjudication is also among the models evaluated on the corrected labels. ``Notes unflagged may be wrong'' means the paper states that unflagged items can also be mislabelled; it is not applicable to full re-annotation.}
\label{tab:literature}
\end{table}
